\section{Limitations}
\label{sec:limitations}

\paragraph{Small sample.}
We use 20 tasks and 3 seeds, so 60 evaluation runs per condition and only 15 on held-out tasks.
Differences of a few tasks are within noise, and some effects we see, such as C8's held-out gain on
qwen, are not significant. More seeds and more tasks would make the results more certain.

\paragraph{Two models.}
We test one weak and one strong small model, and we ran the full set of single-technique conditions
only on the weak one. Other model families, and sizes between 1B and 5B, may behave differently.

\paragraph{Short and simple tasks.}
Our tasks take a few steps each. Continual Harness was designed for long runs of thousands of steps,
where there is much more experience to learn from. With 15 train tasks, the refiner runs only 6 to 18
times per condition. For gemma, the baseline is close to 100\%, which leaves little room for gains.

\paragraph{One version of each technique.}
Each technique has settings we did not tune, such as the number of gate reruns (3), the number of
retrieved notes (5) and the refiner model. We tested each technique alone and all together, but not
all five together minus one (a leave-one-out study), so we cannot say how much each technique adds
inside the full system.

\paragraph{A re-implementation.}
We re-implemented Continual Harness for local models and our tasks, keeping its design as close as
we could. Some changes were needed (Section~\ref{sec:refiner}), so our C1 is a faithful port, not
the original code.

\paragraph{Different servers for the two models.}
All conditions of one model ran on the same server, so comparisons within a model are fair. The two
models ran on different servers and Ollama versions, so exact numbers should not be compared across
models, only the direction of the effects.
